%
% 13-homotopie.tex -- Homotopie
%
% (c) 2026 Prof Dr Andreas Müller
%
\subsection{Homotopie
\label{buch:integration:wegintegral:subsection:homotopie}}
Im vorangegangenen Abschnitt wurde klar, dass das unter geeigneten
Voraussetzungen an das Gebiet das Wegintegral für verschiedene Wege 
zwischen $a$ und $b$ gleich ist.
Die Bedingung war etwas schwerfällig zu formulieren, ergab sich aber
unmittelbar aus dem Satz von Green.
Es musste verlangt werden, zwei Wege zwischen $a$ und $b$ ein in
$U$ enthaltenes Teilgebiet beranden.

Eine etwas intuitivere Formulierung wäre, dass sich die beiden
Wege ineinander deformieren lassen.
Während der Deformation wird das Teilgebiet überstrichen.
\index{Deformation}%
Wir versuchen jetzt, dieses Idee mit einer mathematisch präzisen
Definition zu erfassen.
\input{chapters/030-integration/fig/fig-homotopie.tex}

\begin{definition}[Homotopie von Wegen]
\label{buch:integration:definition:homotpie}
Eine \emph{Homotopie} von Wegen $\gamma_0\colon I\to U$
\index{Homotpie}%
und
$\gamma_1\colon I\to U$, $I=[0,1]$,
in einem Gebiet $U$
mit gleichen Endpunkten
$\gamma_0(0)=\gamma_1(0)$
und
$\gamma_0(1)=\gamma_1(1)$
ist eine Abbildung
\[
H
\colon
I\times I \to U
:
(s,t) \mapsto H(s,t)
\]
mit $H(0,t)=\gamma_0(t)$ und $H(1,t)=\gamma_1(t)$
und konstanten Endpunkten
$H(s,0)=\gamma_i(0)$
und
$H(s,1)=\gamma_i(1)$.
Zwei Wege $\gamma_i\colon I\to U$  heissen \emph{homotop}, in Zeichen
\index{homotop}
$\gamma_0\simeq \gamma_1$, wenn es eine
Homotopie gibt, die die beiden Wege verbindet.
\end{definition}

\begin{satz}
\label{buch:integration:satz:homotopie}
Sei $f\colon U\to\mathbb{C}$ eine holomorphe Funktion und
\[
H
\colon
I\times I
\to
U
\]
eine Homotopie der Wege $\gamma_i$, $i=0,1$, mit $H(s,0) = a$ und $H(s,1)=b$
für $s\in I$.
Dann ist
\[
\int_{\gamma_0} f(z)\,dz
=
\int_{\gamma_1} f(z)\,dz.
\]
\end{satz}

\begin{proof}
Da die Homotopie $H$ eine Abbildung nach $U$ ist, verlaufen
alle Wege
\[
\gamma_s
\colon
I\to\mathbb{C}
:
t
\mapsto
\gamma_s(t) = H(t,s)
\]
in $U$.
Die Vereinigung aller Wege
\[
\bigcup_{0\le s\le 1} \operatorname{im}\gamma_s
=
\{
H(t,s)
\mid
t,s\in I
\}
\subset U
\]
ist in $U$ enhalten, so dass zwei beleibige Kurven immer eine
offene Teilmenge $U$ beranden.
Der Satz~\ref{buch:integration:wegintegral:satz:stammfunktion}
bedeutet daher, dass alle Wegintegrale entlange $\gamma_s$
übereinstimmen, insbesondere auch
\[
\int_{\gamma_0} f(z) \, dz
=
\int_{\gamma_s} f(z) \, dz
= 
\int_{\gamma_1} f(z) \, dz.
\]
Damit ist der Satz bewiesen.
\end{proof}

Der Satz~\ref{buch:integration:satz:homotopie}
sagt, dass ineinander deformierbare Wege zum gleichen Wegintegral führen.
Die Bedingung ist jedoch nicht immer einfach zu verifizieren.
Einfacher wird es, wenn ein Weg zusammenziehbar ist im Sinne
der folgenden Definition.

\begin{definition}[zusammenziehbar]
\label{buch:integration:wegintegral:definition:zusammenziehbar}
Ein geschlossener Weg $\gamma$ in einem Gebiet $U$ heisst
\emph{zusammenziehbar}, wenn der homotop ist zum konstanten Weg
\index{zusammenziehbar}%
$t\mapsto \gamma_1(t)=\gamma(0)=\gamma(1)$.
\end{definition}

In $\mathbb{C}$ ist jeder Weg zusammenziehbar, indem man die Homotopie
\[
H(t,s)
=
(1-s)
\gamma(0)
+
s
\gamma(t) 
\]
verwendet.
Für $s=0$ ist $H(t,0) = \gamma(0)$ der konstante Weg im Punkt
$\gamma(0)$, währen für $s=1$
$H(t,1)=\gamma(t)$ der Weg $\gamma$ ist.
Ein Weg ist nicht zusammenziehbar, wenn er ein ``Loch'' in
der Menge umläuft, da er über das Loch hinweggezogen werden
müsste.
Zum Beispiel ist die Funktion $f(z)=\frac1z$ holomorph in
der Menge $U=\mathbb{C}\setminus\{0\}$ definiert, die an der Stelle
$0$ ein Loch hat.
Ein Weg, der den Nullpunkt umläuft, kann nicht in einen Punkt
zusammengezogen werden, da er dabei den Nullpunkt überqueren müsste.

Mit dem Begriff des zusammenziehbaren Weges ist jetzt
der Satz~\ref{buch:integration:satz:homotopie}
sehr viel anschaulicher zu formulieren.
\input{chapters/030-integration/fig/fig-deformation.tex}

\begin{satz}[Cauchy]
\index{Satz!Cauchy-Integralsatz}%
\index{Cauchy-Integralsatz}
\label{buch:integration:wegintegral:satz:zusammenziehbar}
Ist $\gamma$ ein zusammenziehbarer Weg im Definitionsgebiet einer
holomorphen Funktion, dann gilt
\[
\oint_\gamma f(z)\, dz = 0.
\]
\end{satz}

\begin{proof}
Sei $H$ ein Homotopie zwischen dem konstanten Weg im Punkt $a$ und
der geschlossenen Kurve $\gamma$, also
\[
H(t,0) = a
\qquad
\text{und}
\qquad
H(t,1) = \gamma(t).
\]
Ausserdem soll jeder Weg $t\mapsto H(t,s)$ ein geschlossener Pfad im
Punkt $a$ sein, also $H(0,s)=H(1,s)=a$.
Abbildung~\ref{buch:integration:wegintegral:fig:deformation}
ist die Kurve $\gamma$ rot eingezeichnet.

Der Punkt $b=H(\frac12,1)$ teilt die Kurve $\gamma$ in zwei Kurven
$\gamma_1$ und $\gamma_2$ von $a$ nach $b$ auf.
Die Abbildung $s\mapsto H(\frac12,s)$ ist ein Weg von $a$ nach
$b=H(\frac12,1)$, der in
Abbildung~\ref{buch:integration:wegintegral:fig:deformation}
grün eingezeichnet ist.
Wir wollen zeigen, dass beide Wege $\gamma_1$ und $\gamma_2$ sich in
die grüne Kurve $s\mapsto H(\frac12,s)$ deformieren lässt.

Die Funktion 
\[
K_1
\colon
I\times I
\to
\mathbb{C}
:
(t,s)
\mapsto
K_1(t,s)
=
\begin{cases}
H(t/2s,s)    &\qquad\text{für $t  <  s$}
\\
H(\frac12, t)&\qquad\text{für $s \ge t$}
\end{cases}
\]
deformiert den Weg $\gamma_1$ in den grünen Weg,
denn für $s=1$ entsteht der Weg $\gamma_1$ und
für $s=0$ der grüne Weg.
Für $s$-Werte zwischen $0$ und $1$ folgt der Weg zuerst
der Kurve $t\mapsto H(t,s)$ bis zur grünen Kurve, dann der grünen
Kurve.

Eine ähnliche Deformation des Weges $\gamma_2$ in die grüne
Kurve wird durch
\[
K_2
\colon
I\times I
\to
\mathbb{C}
:
(t,s)
\mapsto
K_2(t,s)
=
\begin{cases}
H((1-2t)/2s,s)    &\qquad\text{für $t  <   s$}
\\
H(\frac12, t)     &\qquad\text{für $s \ge t$}
\end{cases}
\]
gegeben.

Da beide Wege $\gamma_1$ und $\gamma_2$ homotop sind zur
grünen Kurve, sind auch die Wegintegrale entlang der Kurven $\gamma_1$
und $\gamma_2$ gleich dem Wegintegral entlang der grünen Kurve.
Das Wegintegral entlang der Kurve $\gamma$ verschwindet daher.
\end{proof}


